\subsection{Vectorial expression of the convolution product}

PROPOSITION 6. --- Let X,Y,Z be locally compact spaces, \(\varphi\) a continuous mapping of \(X \times Y\) into Z, and \(\lambda, \mu\) measures on X,Y. For \(\lambda\) and \(\mu\) to be \(\varphi\) -convolvable, it is necessary and sufficient that the mapping \((x,y) \mapsto \varepsilon_{\varphi(x,y)} = \varepsilon_x * \varepsilon_y\) of \(X \times Y\) into \(\mathcal{M}(Z)\) be scalarly \((\lambda \otimes \mu)\) -integrable for the topology \(\sigma(\mathcal{M}(Z), \mathcal{K}(Z))\) , in which case

\[
\lambda * \mu = \int_ {\mathrm{X} \times \mathrm{Y}} \left(\varepsilon_ {x} * \varepsilon_ {y}\right) d \lambda (x) d \mu (y).
\]

To say that \(\lambda\) and \(\mu\) are \(\varphi\) -convolvable signifies that, for every \(f \in \mathcal{H}(\mathbf{Z})\) , \(f \circ \varphi\) is \((\lambda \otimes \mu)\) -integrable, that is, for every \(f \in \mathcal{H}(\mathbf{Z})\) the function \((x,y) \mapsto \langle f, \varepsilon_{\varphi(x,y)} \rangle\) is \((\lambda \otimes \mu)\) -integrable, that is, again, that the mapping \((x,y) \mapsto \varepsilon_{\varphi(x,y)}\) of \(\mathrm{X} \times \mathrm{Y}\) into \(\mathcal{M}(\mathbf{Z})\) is scalarly \((\lambda \otimes \mu)\) -integrable for \(\sigma (\mathcal{M}(\mathbb{Z}),\mathcal{K}(\mathbb{Z}))\) . If this is the case, then

\[
\langle \lambda * \mu , f \rangle = \int f (\varphi (x, y)) d \lambda (x) d \mu (y) = \int_ {\mathrm{X} \times \mathrm{Y}} \left\langle \varepsilon_ {\varphi (x, y)}, f \right\rangle d \lambda (x) d \mu (y),
\]

whence \(\lambda * \mu = \int_{\mathbf{X} \times \mathbf{Y}} \varepsilon_{\varphi(x, y)} d\lambda(x) d\mu(y)\) .

PROPOSITION 7. --- Let X,Y,Z be locally compact spaces, \(\varphi\) a continuous mapping of \(X \times Y\) into Z, and \(\lambda, \mu\) measures on X,Y. Assume that for every \(x \in X\) , \(\varepsilon_x\) and \(\mu\) are \(\varphi\) -convolvable. For \(\lambda\) and \(\mu\) to be \(\varphi\) -convolvable, it is necessary and sufficient that the mapping \(x \mapsto \varepsilon_x * |\mu|\) of X into \(\mathcal{M}(Z)\) be scalarly \(\lambda\) -integrable for the topology \(\sigma(\mathcal{M}(Z), \mathcal{K}(Z))\) , in which case \(\lambda * \mu = \int_{X} (\varepsilon_x * \mu) d\lambda(x)\) .

Suppose that \(\lambda\) and \(\mu\) are \(\varphi\) -convolvable. For every \(f \in \mathcal{K}(\mathbb{Z})\) , \(f \circ \varphi\) is \((|\lambda| \otimes |\mu|)\) -integrable, therefore the function \(x \mapsto \int_{\mathbb{Y}} f(\varphi(x, y)) d|\mu|(y) = \langle f, \varepsilon_x * |\mu| \rangle\) (which by hypothesis is defined for all \(x \in \mathbb{X}\) ) is \(\lambda\) -integrable; thus \(x \mapsto \varepsilon_x * |\mu|\) is scalarly \(\lambda\) -integrable for \(\sigma(\mathcal{M}(\mathbb{Z}), \mathcal{K}(\mathbb{Z}))\) , and

\[
\langle f, \lambda * \mu \rangle = \int_ {\mathrm{X}} d \lambda (x) \int_ {\mathrm{Y}} f (\varphi (x, y)) d \mu (y) = \int_ {\mathrm{X}} \left\langle f, \varepsilon_ {x} * \mu \right\rangle d \lambda (x),
\]

whence \(\lambda * \mu = \int_{\mathbf{X}} (\varepsilon_x * \mu) d\lambda(x)\) . Conversely, suppose that the mapping \(x \mapsto \varepsilon_x * |\mu|\) of X into \(\mathcal{M}(\mathbf{Z})\) is scalarly \(\lambda\) -integrable for \(\sigma(\mathcal{M}(\mathbf{Z}), \mathcal{K}(\mathbf{Z}))\) . Let \(f \in \mathcal{H}_+(\mathbf{Z})\) . Then the function \((x,y) \mapsto f(\varphi(x,y))\) is continuous and (Ch. V, §8, No.~3, Prop. 5)

\[
\begin{array}{r l} \iint^ {*} f (\varphi (x, y)) d | \lambda | (x) d | \mu | (y) & = \int^ {*} d | \lambda | (x) \int^ {*} f (\varphi (x, y)) d | \mu | (y) \\ & = \int^ {*} \langle f, \varepsilon_ {x} * | \mu | \rangle d | \lambda | (x) <   + \infty . \end{array}
\]

Therefore \(f \circ \varphi\) is \((\lambda \otimes \mu)\) -integrable, so that \(\lambda\) and \(\mu\) are \(\varphi\) -convolvable.

